\begin{proofbox}
\textbf{量词否定的逻辑等价}（以全称量词为例）：
\[
\lnot\big(\forall x\in U,\;P(x)\big) \;\equiv\; \exists x\in U,\;\lnot P(x).
\]
\textbf{证明}：若"所有 $x$ 都满足 $P$"为假，则至少存在一个 $x$ 使 $P$ 不成立，反之亦然。这是逻辑基本公理。
\end{proofbox}

\begin{proofbox}
\textbf{德摩根定律的证明}（以 $p\land q$ 为例）：
\begin{center}
\begin{tabular}{|c|c|c|c|c|}
\hline
$p$ & $q$ & $p\land q$ & $\lnot(p\land q)$ & $(\lnot p)\lor(\lnot q)$ \\ \hline
真 & 真 & 真 & 假 & 假 \\
真 & 假 & 假 & 真 & 真 \\
假 & 真 & 假 & 真 & 真 \\
假 & 假 & 假 & 真 & 真 \\ \hline
\end{tabular}
\end{center}
两列真值完全相同，故 $\lnot(p\land q) \equiv (\lnot p)\lor(\lnot q)$。类似可证另一个。
\end{proofbox}

\begin{proofbox}
\textbf{比较关系的否定}（以"大于"为例）：
\[
\lnot(a>b) \equiv a\le b.
\]
实数的序关系具有三歧性：对任意 $a,b$，$a>b$、$a=b$、$a<b$ 有且仅有一个成立。因此"$a>b$"为假等价于"$a\le b$"为真。
\end{proofbox}